\section{核验锚点簇：洪熙至天顺：边镇、土木与复辟后的重组（1425--1464）}
\label{register:1425-1464_hongxi_tianshun}
这一时期的危机既发生在北方行军与北京守备之间，也发生在继承程序和财政调度之中。每条首次描写仅确认账本能支持的行动链，并把数字、路线和动机留在可继续核对的位置。

以下锚点只承担年序导航；本章前文负责叙事。完整来源、校勘与材料限度留在来源附录和必要的信息框，不把账本字段重复铺成读者段落。

\begin{description}
  \item[\eventanchor{EM-1425-XUANDE-ACCESSION}{1425（洪熙元年）五月：仁宗去世，皇太子朱瞻基即位为宣德帝}]
  \item[\eventanchor{ML-1425-001}{1425（洪熙元年）五月29日：洪熙帝驾崩}]
  \item[\eventanchor{ML-1425-002}{1425（洪熙元年）六月27日：朱瞻基即位宣德皇帝}]
  \item[\eventanchor{ML-1425-003}{1425（洪熙元年）九月2日：朱高煦造反}]
  \item[\eventanchor{ML-1425-004}{1425（洪熙元年）九月22日：朱高煦战败}]
  \item[\eventanchor{EM-1426-GAOXU-REBELLION}{1426（宣德元年）：汉王朱高煦在乐安起兵，宣德帝亲征后其降伏并被处置}]
  \item[\eventanchor{ML-1426-001}{1426（宣德元年）：明朝派遣一失哈到野女真建造船坞和仓库}]
  \item[\eventanchor{ML-1426-002}{1426（宣德元年）十月5日：林山起义：黎来的军队在二同–秋同战役中对明朝的进攻造成重大伤亡}]
  \item[\eventanchor{ML-1426-003}{1426（宣德元年）冬：林山起义：林山军队将明军驱逐出红河三角洲和越南北部的大部分地区}]
  \item[\eventanchor{ML-1427-001}{1427（宣德二年）十月10日：林山起义：明军援军在林山被围歼}]
  \item[\eventanchor{ML-1427-002}{1427（宣德二年）十二月14日：林山起义：明军撤出交趾}]
  \item[\eventanchor{ML-1427-003}{1427（宣德二年）：本年朝廷围绕继承、官僚协调和政令传达形成连续记录，可用于核对宣德至正统的中枢运作。}]
  \item[\eventanchor{EM-1428-LE-LOI-RECOGNITION}{1428（宣德三年）：明接受黎利政权的朝贡关系并调整安南外交称谓}]
  \item[\eventanchor{ML-1428-001}{1428（宣德三年）三月25日：宝藏远航：宣德皇帝命郑和督修大报恩寺}]
  \item[\eventanchor{ML-1428-002}{1428（宣德三年）四月29日：黎黎在后黎王朝的领导下重建了大越王国}]
  \item[\eventanchor{ML-1428-003}{1428（宣德三年）十月：乌连开袭明边境，宣德皇帝亲自率兵击退}]
  \item[\eventanchor{ML-1429-001}{1429（宣德四年）：宣德皇帝在北京郊外进行大检阅}]
  \item[\eventanchor{ML-1429-002}{1429（宣德四年）：明军使用携带手炮的骑兵。}]
  \item[\eventanchor{ML-1429-003}{1429（宣德四年）：漕运、粮饷与军户事务的年度文书显示财政—军事相互约束，定额不能替代实际征解。}]
  \item[\eventanchor{ML-1429-004}{1429（宣德四年）：交趾、北边或辽东的军令与边报在本年构成军事决策线索，报称战果应与对方及地方材料比照。}]
  \item[\eventanchor{EM-1430-YUQIAN-INSPECTION}{1430（宣德五年）：于谦奉命巡抚河南、山西，处理灾荒、军政与地方事务}]
  \item[\eventanchor{EM-1430-ZHENGHE-SEVENTH-VOYAGE}{1430（宣德五年）：郑和第七次下西洋自南京出发}]
  \item[\eventanchor{ML-1430-001}{1430（宣德五年）五月：宣德皇帝下令减免所有皇家土地的赋税}]
  \item[\eventanchor{ML-1430-002}{1430（宣德五年）六月29日：宝藏下西洋：宣德皇帝下令第七次下西洋}]
  \item[\eventanchor{EM-1431-TIANFEI-STELE}{1431（宣德六年）：郑和等立《天妃灵应之记》碑，记述历次航行与祈祷}]
  \item[\eventanchor{ML-1431-001}{1431（宣德六年）一月19日：宝藏之旅：宝藏船队从南京出发}]
  \item[\eventanchor{ML-1431-002}{1431（宣德六年）三月14日：宝藏之旅：刘家港碑刻立碑}]
  \item[\eventanchor{ML-1431-003}{1431（宣德六年）六月12日：越南黎朝皇帝黎太祖与明朝中国建立朝贡关系，并被明朝皇帝封为安南王。}]
  \item[\eventanchor{ML-1432-001}{1432（宣德七年）九月12日：宝藏之旅：宝藏船队抵达萨穆德拉巴赛苏丹国，红宝、马欢随船队访问孟加拉}]
  \item[\eventanchor{ML-1432-002}{1432（宣德七年）：明朝派遣一失哈向明朝联军女真人进贡印玺并修复永宁寺}]
  \item[\eventanchor{ML-1432-003}{1432（宣德七年）：朝贡、册封和边地交涉的记录为本年多政权关系留档，礼仪语言与实际控制范围必须区分。}]
  \item[\eventanchor{ML-1432-004}{1432（宣德七年）：漕运、粮饷与军户事务的年度文书显示财政—军事相互约束，定额不能替代实际征解。}]
  \item[\eventanchor{EM-1433-ZHENGHE-RETURN}{1433（宣德八年）：第七次航行舰队返抵中国，郑和卒年与卒地存在不同说法}]
  \item[\eventanchor{ML-1433-001}{1433（宣德八年）：宝藏航行：郑和去世}]
  \item[\eventanchor{ML-1433-002}{1433（宣德八年）：宝藏之旅：洪宝和马欢抵达卡利卡特并派七人前往麦加，洪宝访问了乔法尔、拉萨、亚丁、摩加迪沙和巴拉瓦，然后返回中国}]
  \item[\eventanchor{ML-1433-003}{1433（宣德八年）三月9日：宝藏之旅：宝船队从霍尔木兹出发返回中国}]
  \item[\eventanchor{ML-1433-004}{1433（宣德八年）六月：日本使团访问明朝中国：明日关系恢复}]
  \item[\eventanchor{ML-1433-005}{1433（宣德八年）七月7日：宝藏之旅：宝船队回到中国}]
  \item[\eventanchor{ML-1433-006}{1433（宣德八年）九月14日：宝藏航行：萨穆德拉·帕赛苏丹国、卡利卡特、科钦、锡兰、乔法尔、亚丁、哥印拜陀、霍尔木兹、卡亚尔和麦加的使节进贡}]
  \item[\eventanchor{ML-1433-007}{1433（宣德八年）：宝藏之旅：马欢出版《莺芽圣澜》}]
  \item[\eventanchor{ML-1434-001}{1434（宣德九年）：宝藏之旅：龚震出版《西洋繁国志》}]
  \item[\eventanchor{ML-1434-002}{1434（宣德九年）：本年朝廷围绕继承、官僚协调和政令传达形成连续记录，可用于核对宣德至正统的中枢运作。}]
  \item[\eventanchor{ML-1434-003}{1434（宣德九年）：学校、科举和礼制的本年政策材料为社会流动和国家知识秩序提供可检索锚点。}]
  \item[\eventanchor{ML-1434-004}{1434（宣德九年）：朝贡、册封和边地交涉的记录为本年多政权关系留档，礼仪语言与实际控制范围必须区分。}]
  \item[\eventanchor{EM-1435-XUANDE-DEATH}{1435（宣德十年）正月：宣德帝去世，九岁皇太子朱祁镇即位为英宗}]
  \item[\eventanchor{ML-1435-001}{1435（宣德十年）一月31日：宣德皇帝驾崩，张皇后（洪熙）摄政为正统皇帝}]
  \item[\eventanchor{ML-1435-002}{1435（宣德十年）：华北平原及山东遭遇旱灾、瘟疫}]
  \item[\eventanchor{EM-1436-LUCHUAN-CRISIS}{1436（正统元年）：麓川思任发扩张引发边境冲突，明廷开始部署征讨}]
  \item[\eventanchor{ML-1436-001}{1436（正统元年）：宝藏航行：明朝禁止建造海船}]
  \item[\eventanchor{ML-1436-002}{1436（正统元年）：宝藏之旅：费辛出版《星茶圣澜》}]
  \item[\eventanchor{ML-1436-003}{1436（正统元年）：苏北、华北平原、山东等地遭受洪涝灾害}]
  \item[\eventanchor{ML-1437-001}{1437（正统二年）：山西、陕西遭遇干旱}]
  \item[\eventanchor{ML-1437-002}{1437（正统二年）：苏北地区发生洪涝灾害}]
  \item[\eventanchor{ML-1437-003}{1437（正统二年）：灾情上报、赈恤或河工材料标记本年地方风险，损失与救济效果需回到受灾地区核实。（材料核查重点：诏令或奏议的形成与发受链）}]
  \item[\eventanchor{ML-1437-004}{1437（正统二年）：灾情上报、赈恤或河工材料标记本年地方风险，损失与救济效果需回到受灾地区核实。（材料核查重点：报称信息的来源、抄传与行政分类）}]
  \item[\eventanchor{EM-1438-FIRST-LUCHUAN-CAMPAIGN}{1438（正统三年）：王骥等统兵发动第一次大规模征麓川行动}]
  \item[\eventanchor{ML-1438-001}{1438（正统三年）十二月8日：陆川平绵之战：明讨伐勐茂司仁法，攻打邻近土司，但未能击败他}]
  \item[\eventanchor{ML-1438-002}{1438（正统三年）：本年朝廷围绕继承、官僚协调和政令传达形成连续记录，可用于核对宣德至正统的中枢运作。（材料核查重点：诏令或奏议的形成与发受链）}]
  \item[\eventanchor{ML-1438-003}{1438（正统三年）：本年朝廷围绕继承、官僚协调和政令传达形成连续记录，可用于核对宣德至正统的中枢运作。（材料核查重点：报称信息的来源、抄传与行政分类）}]
  \item[\eventanchor{ML-1439-001}{1439（正统四年）：华北平原及山东地区发生洪涝灾害}]
  \item[\eventanchor{ML-1439-002}{1439（正统四年）：交趾、北边或辽东的军令与边报在本年构成军事决策线索，报称战果应与对方及地方材料比照。（材料核查重点：诏令或奏议的形成与发受链）}]
  \item[\eventanchor{ML-1439-003}{1439（正统四年）：交趾、北边或辽东的军令与边报在本年构成军事决策线索，报称战果应与对方及地方材料比照。（材料核查重点：报称信息的来源、抄传与行政分类）}]
  \item[\eventanchor{ML-1439-004}{1439（正统四年）：交趾、北边或辽东的军令与边报在本年构成军事决策线索，报称战果应与对方及地方材料比照。（材料核查重点：同年地方材料所见的执行范围）}]
  \item[\eventanchor{EM-1440-ESEN-RAID}{1440（正统五年）：也先势力进犯明北边，边镇紧张加剧}]
  \item[\eventanchor{ML-1440-001}{1440（正统五年）：苏州、江南、华北平原、山东等地遭受洪涝灾害}]
  \item[\eventanchor{ML-1440-002}{1440（正统五年）：浙江饥荒袭来}]
  \item[\eventanchor{EM-1441-SECOND-LUCHUAN-CAMPAIGN}{1441（正统六年）：明廷再遣大军征麓川，试图追捕思任发并重建土司秩序}]
  \item[\eventanchor{ML-1441-001}{1441（正统六年）二月27日：陆川-平绵之战：明军进攻蒙茂}]
  \item[\eventanchor{ML-1441-002}{1441（正统六年）：华北平原及山东地区发生洪涝灾害}]
  \item[\eventanchor{ML-1441-003}{1441（正统六年）：浙江饥荒袭来}]
  \item[\eventanchor{EM-1442-ZHANG-TAIHOU-DEATH}{1442（正统七年）：张太皇太后去世，幼主时期的辅政平衡发生变化}]
  \item[\eventanchor{ML-1442-001}{1442（正统七年）一月：陆川平绵之战：蒙茂战败，司仁法逃往阿瓦}]
  \item[\eventanchor{ML-1442-002}{1442（正统七年）十一月20日：张皇后（洪熙）去世}]
  \item[\eventanchor{ML-1442-003}{1442（正统七年）：朝贡、册封和边地交涉的记录为本年多政权关系留档，礼仪语言与实际控制范围必须区分。}]
  \item[\eventanchor{EM-1443-THIRD-LUCHUAN-CAMPAIGN}{1443（正统八年）：明廷继续第三次征麓川并以招抚、分置应对残余势力}]
  \item[\eventanchor{ML-1443-001}{1443（正统八年）三月：陆川平绵之战：明军击败司继发，但未能俘获他}]
  \item[\eventanchor{ML-1443-002}{1443（正统八年）：漕运、粮饷与军户事务的年度文书显示财政—军事相互约束，定额不能替代实际征解。}]
  \item[\eventanchor{ML-1443-003}{1443（正统八年）：学校、科举和礼制的本年政策材料为社会流动和国家知识秩序提供可检索锚点。}]
  \item[\eventanchor{ML-1444-001}{1444（正统九年）：山西、陕西发生饥荒}]
  \item[\eventanchor{ML-1444-002}{1444（正统九年）：苏北地区发生洪涝灾害}]
  \item[\eventanchor{ML-1444-003}{1444（正统九年）：朝贡、册封和边地交涉的记录为本年多政权关系留档，礼仪语言与实际控制范围必须区分。}]
  \item[\eventanchor{ML-1444-004}{1444（正统九年）：灾情上报、赈恤或河工材料标记本年地方风险，损失与救济效果需回到受灾地区核实。}]
  \item[\eventanchor{EM-1445-ESEN-TRIBUTE-CONFLICT}{1445（正统十年）：也先借朝贡、互市和使团规模问题向明廷施压，北边交涉恶化}]
  \item[\eventanchor{ML-1445-001}{1445（正统十年）八月：陆川-平绵战役：阿瓦将斯仁法交给明朝，以换取明朝支持进攻鲜威}]
  \item[\eventanchor{ML-1445-002}{1445（正统十年）：浙江干旱鼠疫肆虐}]
  \item[\eventanchor{ML-1446-002}{1446（正统十一年）：江南洪水泛滥}]
  \item[\eventanchor{ML-1446-003}{1446（正统十一年）：学校、科举和礼制的本年政策材料为社会流动和国家知识秩序提供可检索锚点。}]
  \item[\eventanchor{ML-1447-001}{1447（正统十二年）：叶宗柳与浙江银矿工人叛乱}]
  \item[\eventanchor{ML-1447-002}{1447（正统十二年）：苏北饥荒袭来}]
  \item[\eventanchor{ML-1447-003}{1447（正统十二年）：灾情上报、赈恤或河工材料标记本年地方风险，损失与救济效果需回到受灾地区核实。}]
  \item[\eventanchor{EM-1448-DENG-MAOQI}{1448（正统十三年）：福建邓茂七起事扩大，地方官军与中央调兵镇压}]
  \item[\eventanchor{EM-1448-THIRD-LUCHUAN-EXPEDITION}{1448（正统十三年）：明廷为麓川问题再度大规模调兵，征发延及多省军镇}]
  \item[\eventanchor{ML-1448-001}{1448（正统十三年）三月：邓茂奇与一群佃农在闽赣边境西北叛乱}]
  \item[\eventanchor{ML-1448-002}{1448（正统十三年）十二月：明军杀了叶宗六，但他的叛军毫发无伤，继续向南撤退，围困滁州}]
  \item[\eventanchor{ML-1448-003}{1448（正统十三年）：1448年黄河洪水：黄河堤决口}]
  \item[\eventanchor{ML-1448-004}{1448（正统十三年）：干旱和蝗灾袭击中国西北地区}]
  \item[\eventanchor{ML-1448-005}{1448（正统十三年）：江南大旱}]
  \item[\eventanchor{EM-1449-BEIJING-DEFENSE}{1449（正统十四年）九月：于谦等组织北京守备，拒绝南迁并抵御瓦剌进逼}]
  \item[\eventanchor{EM-1449-JINGTAI-ACCESSION}{1449（正统十四年）九月：郕王朱祁钰即帝位为景泰帝，尊英宗为太上皇}]
  \item[\eventanchor{ML-1449-001}{1449（正统十四年）三月：陆川平绵之战：明军因藏匿司继发而入侵蒙阳，司继发又逃脱}]
  \item[\eventanchor{ML-1449-002}{1449（正统十四年）五月：邓茂奇叛军被击败}]
  \item[\eventanchor{ML-1449-003}{1449（正统十四年）七月：图木危机：北元事实上的统治者瓦剌也先太师入侵明朝}]
  \item[\eventanchor{ML-1449-004}{1449（正统十四年）八月4日：土木危机：正统皇帝出京亲自迎战也先太史}]
  \item[\eventanchor{ML-1449-006}{1449（正统十四年）八月：叶宗流的叛军被击败}]
  \item[\eventanchor{ML-1449-007}{1449（正统十四年）九月1日：土木危机：明军全军覆没，正统皇帝被也先太师俘虏}]
  \item[\eventanchor{ML-1449-008}{1449（正统十四年）九月23日：朱祁钰即位景泰皇帝}]
  \item[\eventanchor{ML-1449-009}{1449（正统十四年）十月27日：也先太史围攻北京未能攻下，五天后撤军}]
  \item[\eventanchor{ML-1449-010}{1449（正统十四年）：黄河堤坝再次决堤导致河道略有改道}]
  \item[\eventanchor{EM-1450-YINGZONG-RETURN}{1450（景泰元年）八月：瓦剌释放英宗，英宗返京后被安置为太上皇}]
  \item[\eventanchor{ML-1450-001}{1450（景泰元年）九月19日：正统皇帝被释放回到北京，被景泰皇帝软禁}]
  \item[\eventanchor{ML-1450-002}{1450（景泰元年）：贵州、湖广瑶族、苗族起义}]
  \item[\eventanchor{ML-1450-003}{1450（景泰元年）：山东发生饥荒}]
  \item[\eventanchor{ML-1450-004}{1450（景泰元年）：军饷、漕运和户籍赋役的本年记录揭示恢复秩序的行政成本，账面额与实收须分列。}]
  \item[\eventanchor{ML-1451-001}{1451（景泰二年）：土木危机后的边镇整顿、京师防务或复辟前后军政调度在本年材料中构成连续问题。（材料核查重点：诏令或奏议的形成与发受链）}]
  \item[\eventanchor{ML-1451-002}{1451（景泰二年）：科举、学校和法律条文在本年制度材料中可追踪，不能据规范文本推定州县实践一致。}]
  \item[\eventanchor{ML-1451-003}{1451（景泰二年）：地方灾异、赈济与水利条目为本年社会史提供入口，地方志的修纂层次需要说明。}]
  \item[\eventanchor{ML-1451-004}{1451（景泰二年）：土木危机后的边镇整顿、京师防务或复辟前后军政调度在本年材料中构成连续问题。（材料核查重点：报称信息的来源、抄传与行政分类）}]
  \item[\eventanchor{ML-1451-005}{1451（景泰二年）：蒙古诸部、朝鲜及西北绿洲的交涉在本年留下多方材料，应以具体使团和地名校正概称。}]
  \item[\eventanchor{ML-1452-001}{1452（景泰三年）：1452年黄河洪水：黄河堤决口}]
  \item[\eventanchor{ML-1452-002}{1452（景泰三年）：瑶族、苗族叛乱被镇压}]
  \item[\eventanchor{ML-1452-003}{1452（景泰三年）：中国北方遭遇洪涝灾害}]
  \item[\eventanchor{ML-1452-004}{1452（景泰三年）：科举、学校和法律条文在本年制度材料中可追踪，不能据规范文本推定州县实践一致。}]
  \item[\eventanchor{ML-1453-001}{1453（景泰四年）：地方灾异、赈济与水利条目为本年社会史提供入口，地方志的修纂层次需要说明。（材料核查重点：诏令或奏议的形成与发受链）}]
  \item[\eventanchor{ML-1453-002}{1453（景泰四年）：土木危机后的边镇整顿、京师防务或复辟前后军政调度在本年材料中构成连续问题。}]
  \item[\eventanchor{ML-1453-003}{1453（景泰四年）：景泰、天顺时期的继承、人事与诏令链条可在本年材料中定位，政变责任不可由单一史书裁断。}]
  \item[\eventanchor{ML-1453-004}{1453（景泰四年）：地方灾异、赈济与水利条目为本年社会史提供入口，地方志的修纂层次需要说明。（材料核查重点：报称信息的来源、抄传与行政分类）}]
  \item[\eventanchor{MM-1453-ESEN-KHAN}{1453（景泰四年）：也先杀脱脱不花后自称大汗，瓦剌内部权力关系重组}]
  \item[\eventanchor{ML-1454-001}{1454（景泰五年）：异常大雪导致苏州、杭州饥荒}]
  \item[\eventanchor{ML-1454-002}{1454（景泰五年）：土木危机后的边镇整顿、京师防务或复辟前后军政调度在本年材料中构成连续问题。}]
  \item[\eventanchor{ML-1454-003}{1454（景泰五年）：景泰、天顺时期的继承、人事与诏令链条可在本年材料中定位，政变责任不可由单一史书裁断。}]
  \item[\eventanchor{ML-1454-004}{1454（景泰五年）：军饷、漕运和户籍赋役的本年记录揭示恢复秩序的行政成本，账面额与实收须分列。}]
  \item[\eventanchor{ML-1455-001}{1455（景泰六年）：徐有贞完成黄河堤防修缮工作}]
  \item[\eventanchor{ML-1455-002}{1455（景泰六年）：华中地区大面积干旱}]
  \item[\eventanchor{ML-1455-003}{1455（景泰六年）：蒙古诸部、朝鲜及西北绿洲的交涉在本年留下多方材料，应以具体使团和地名校正概称。}]
  \item[\eventanchor{MM-1455-ESEN-DEATH}{1455（景泰六年）：也先在草原内部冲突中被杀，瓦剌汗权继承发生断裂}]
  \item[\eventanchor{ML-1456-001}{1456（景泰七年）：湖广苗族叛乱遭镇压}]
  \item[\eventanchor{ML-1456-002}{1456（景泰七年）：景泰、天顺时期的继承、人事与诏令链条可在本年材料中定位，政变责任不可由单一史书裁断。}]
  \item[\eventanchor{ML-1456-003}{1456（景泰七年）：军饷、漕运和户籍赋役的本年记录揭示恢复秩序的行政成本，账面额与实收须分列。}]
  \item[\eventanchor{ML-1456-004}{1456（景泰七年）：科举、学校和法律条文在本年制度材料中可追踪，不能据规范文本推定州县实践一致。}]
  \item[\eventanchor{ML-1457-001}{1457（天顺元年）二月11日：前皇帝被军方复职，成为天顺皇帝}]
  \item[\eventanchor{MM-1457-JINGTAI-DEATH}{1457（天顺元年）正月：英宗复位后景泰帝不久去世，景泰朝的皇位主张随之终结}]
  \item[\eventanchor{MM-1457-YU-QIAN-EXECUTION}{1457（天顺元年）正月：于谦以“谋逆”罪被处死}]
  \item[\eventanchor{ML-1458-001}{1458（天顺二年）：军饷、漕运和户籍赋役的本年记录揭示恢复秩序的行政成本，账面额与实收须分列。}]
  \item[\eventanchor{ML-1458-002}{1458（天顺二年）：蒙古诸部、朝鲜及西北绿洲的交涉在本年留下多方材料，应以具体使团和地名校正概称。}]
  \item[\eventanchor{ML-1458-003}{1458（天顺二年）：地方灾异、赈济与水利条目为本年社会史提供入口，地方志的修纂层次需要说明。}]
  \item[\eventanchor{MM-1458-CROWN-PRINCE-RESTORED}{1458（天顺二年）：英宗复立朱见深为皇太子}]
  \item[\eventanchor{ML-1459-001}{1459（天顺三年）：景泰、天顺时期的继承、人事与诏令链条可在本年材料中定位，政变责任不可由单一史书裁断。}]
  \item[\eventanchor{ML-1459-002}{1459（天顺三年）：蒙古诸部、朝鲜及西北绿洲的交涉在本年留下多方材料，应以具体使团和地名校正概称。}]
  \item[\eventanchor{ML-1459-003}{1459（天顺三年）：科举、学校和法律条文在本年制度材料中可追踪，不能据规范文本推定州县实践一致。}]
  \item[\eventanchor{ML-1459-004}{1459（天顺三年）：土木危机后的边镇整顿、京师防务或复辟前后军政调度在本年材料中构成连续问题。}]
  \item[\eventanchor{ML-1460-001}{1460（天顺四年）：科举、学校和法律条文在本年制度材料中可追踪，不能据规范文本推定州县实践一致。}]
  \item[\eventanchor{ML-1460-002}{1460（天顺四年）：地方灾异、赈济与水利条目为本年社会史提供入口，地方志的修纂层次需要说明。}]
  \item[\eventanchor{ML-1460-003}{1460（天顺四年）：景泰、天顺时期的继承、人事与诏令链条可在本年材料中定位，政变责任不可由单一史书裁断。}]
  \item[\eventanchor{MM-1460-SHI-HENG-FALL}{1460（天顺四年）：石亨被逮下狱，复辟功臣集团失去中枢优势}]
  \item[\eventanchor{ML-1461-001}{1461（天顺五年）八月7日：曹钦之乱：曹钦叛乱并试图攻入北京，但被捕并被迫自杀}]
  \item[\eventanchor{ML-1461-002}{1461（天顺五年）：土木危机后的边镇整顿、京师防务或复辟前后军政调度在本年材料中构成连续问题。}]
  \item[\eventanchor{ML-1461-003}{1461（天顺五年）：军饷、漕运和户籍赋役的本年记录揭示恢复秩序的行政成本，账面额与实收须分列。}]
  \item[\eventanchor{MM-1461-CAO-QIN-REBELLION}{1461（天顺五年）七月：曹钦在北京发动兵变，攻东安门等处后被官军平定}]
  \item[\eventanchor{ML-1462-001}{1462（天顺六年）：军饷、漕运和户籍赋役的本年记录揭示恢复秩序的行政成本，账面额与实收须分列。}]
  \item[\eventanchor{ML-1462-002}{1462（天顺六年）：地方灾异、赈济与水利条目为本年社会史提供入口，地方志的修纂层次需要说明。}]
  \item[\eventanchor{ML-1462-004}{1462（天顺六年）：蒙古诸部、朝鲜及西北绿洲的交涉在本年留下多方材料，应以具体使团和地名校正概称。}]
  \item[\eventanchor{ML-1463-001}{1463（天顺七年）：蒙古诸部、朝鲜及西北绿洲的交涉在本年留下多方材料，应以具体使团和地名校正概称。}]
  \item[\eventanchor{ML-1463-002}{1463（天顺七年）：土木危机后的边镇整顿、京师防务或复辟前后军政调度在本年材料中构成连续问题。}]
  \item[\eventanchor{ML-1463-003}{1463（天顺七年）：军饷、漕运和户籍赋役的本年记录揭示恢复秩序的行政成本，账面额与实收须分列。}]
  \item[\eventanchor{ML-1463-004}{1463（天顺七年）：科举、学校和法律条文在本年制度材料中可追踪，不能据规范文本推定州县实践一致。}]
  \item[\eventanchor{ML-1464-001}{1464（天顺八年）二月23日：天顺帝驾崩，朱千深即位成化帝}]
  \item[\eventanchor{ML-1464-002}{1464（天顺八年）：广西瑶族起义军侯大狗}]
  \item[\eventanchor{ML-1464-003}{1464（天顺八年）：宝藏航行：宝藏航行的文献被刘大侠从兵部档案中移出并销毁，理由是它们是“离奇之事，远非耳目所证的欺骗性夸大”，“三宝出征西海，浪费了数十万的金钱和粮食，而且死亡的人也可以算进”。虽然他带回了许多珍贵的东西，但这对国家有什么好处呢？}]
  \item[\eventanchor{MM-1464-XIANZONG-ACCESSION}{1464（天顺八年）正月：英宗去世，皇太子朱见深即位，是为宪宗}]
\end{description}
